\section{Method: The \method\ Framework}
\label{sec:method}

We organize memory for VLAs along two orthogonal axes: \emph{time
scale} (hot/cold/persistent) and \emph{modality} (video, text,
retrieval-embedding). Section~\ref{sec:method:architecture} gives the
high-level architecture. Sections~\ref{sec:method:lang},
\ref{sec:method:persistent}, and \ref{sec:method:integration} detail
the language memory, the persistent store, and the planner
integration, respectively.

\subsection{Architecture}
\label{sec:method:architecture}

Figure~\ref{fig:architecture} shows the three-tier design.

\begin{itemize}
\item \textbf{Hot memory (seconds)}: a short-term video/sensor buffer
      such as the sliding window of Mem-0 or the sparse-attention
      video encoder of MEM. We intentionally do \emph{not}
      reimplement hot memory; \method\ composes with whichever hot
      memory the underlying VLA provides.
\item \textbf{Cold memory (minutes)}: a semi-structured language
      memory expressed as human-readable text (see
      Section~\ref{sec:method:lang}). The cold memory is updated at
      subtask boundaries by the high-level planner.
\item \textbf{Persistent memory (days)}: a FAISS-style vector store
      of \emph{episode summaries} indexed by text embedding
      (Section~\ref{sec:method:persistent}). New episodes query the
      store on start-up and archive their final summary on
      termination.
\end{itemize}

All three tiers reach the underlying VLA backbone through a single
touchpoint: the planner prompt. The planner consumes the current
cold memory alongside the recent hot memory tokens and any retrieved
persistent facts, and produces (a) the next subtask and (b) an
updated cold memory.

\subsection{Semi-Structured Language Memory}
\label{sec:method:lang}

Our language memory schema is designed to be \emph{fully
round-trippable}: the same object can be rendered as text for LLM
consumption and parsed back into structured form for programmatic
inspection. Each \texttt{Memory} object captures:

{\small
\begin{verbatim}
[TASK]    <goal>
[SESSION] <id>
[SUBTASKS]
  1. [X] <done subtask>
  2. [>] <current subtask>
  3. [ ] <pending subtask>
[CROSS-SESSION]
  - <fact> (source: <session>, conf=<0..1>)
[FAILURES]
  - <error> -> <intervention>
\end{verbatim}
}

The status markers \texttt{[X]/[>]/[\,]/[!]} explicitly encode
subtask state (done, current, pending, failed). Cross-session facts
carry a provenance tag and a confidence score. Failures record both
the error and the intervention taken. This schema is inspired by
MEM's natural-language memory \cite{mem} but extends it with
structural markers that support \emph{programmatic} operations:
diffing, anomaly detection, and rule-based analysis.

Crucially, the schema preserves the \emph{interpretability} lost in
dense-embedding memories (MemoryVLA, ReMem-VLA). Practitioners can
read the cold memory at any point in an episode to diagnose
failures, and downstream tools can reason over it symbolically.

\subsection{Persistent Cross-Session Memory}
\label{sec:method:persistent}

The persistent layer stores an \texttt{EpisodeSummary} object per
completed episode, containing: session id, task, outcome
(success/failure/partial), narrative, and a typed list of facts
(object-location, task-outcome, failure-cause, user-preference).

At episode start, the planner issues a text query---typically the
global task goal---to retrieve the top-$K$ most similar past
summaries. Results are formatted into a \texttt{<cross-session
memory>} block that is concatenated into the planner prompt before
the \texttt{<finished\_subtasks>} section. At episode end, the final
cold memory is converted into a summary and persisted.

\paragraph{Retrieval.}
Summaries are embedded via a pluggable \texttt{Embedder} protocol.
Our initial release ships two implementations:
(i) \texttt{HashingEmbedder}, a zero-dependency deterministic hasher
used for unit tests and minimal deployments;
(ii) \texttt{SentenceTransformerEmbedder}, wrapping
\texttt{sentence-transformers} for semantic quality.
Replacing the embedder with a Qwen3-VL text encoder is a single
class substitution. The store returns top-$K$ hits by cosine
similarity; we rely on numpy for indexing at $N<50{,}000$
summaries and recommend FAISS for larger scales.

\paragraph{Update protocol.}
A rule-based summarizer converts the final \texttt{Memory} of an
episode into an \texttt{EpisodeSummary}. Completed subtasks become
narrative bullets; failures become \texttt{failure\_cause} facts;
task outcome is inferred from subtask statuses. An LLM-based
summarizer (called once per episode, not per step) is the planned
v0.3 upgrade; the rule-based variant suffices to validate the
persistence layer.

\subsection{Integration with Existing Memory-Aware VLAs}
\label{sec:method:integration}

\method\ is designed for \emph{incremental adoption}. Our primary
integration target is Mem-0 \cite{rmbench}, chosen because its
planner-executor hierarchy already mirrors our cold/hot distinction.
Two edits suffice:

\begin{enumerate}
\item \texttt{MemoryMattersPlanner.\_\_init\_\_} gains a
      \texttt{self.persistent\_store = MemoryStore(...)} attribute.
\item \texttt{MemoryMattersPlanner.prepare\_qwen\_input} calls
      \texttt{inject\_into\_prompt(...)} before returning, which
      inserts a \texttt{<cross-session memory>} block into the Qwen
      prompt.
\end{enumerate}

A third one-line hook in \texttt{MemoryMattersAgent.reset} invokes
\texttt{summarize\_rule\_based} at episode end to persist the
summary. No retraining is required for the integration to take
effect: Qwen3-VL treats the injected block as additional context.

The same pattern applies to MemER \cite{memer} and MemoryVLA
\cite{memoryvla}, with minor adaptations documented in the
release notes. For purely end-to-end VLAs (e.g., CronusVLA), the
persistent memory can be surfaced through a separate high-level
planner layered on top.

\subsection{Training Considerations}
\label{sec:method:training}

The current release operates at inference time only. For training,
we propose a joint objective combining action imitation, memory
summary generation, and consistency between the two:

\begin{equation}
\mathcal{L} = \lambda_a \mathcal{L}_{\text{action}}
            + \lambda_s \mathcal{L}_{\text{subtask}}
            + \lambda_m \mathcal{L}_{\text{memory}}
            + \lambda_c \mathcal{L}_{\text{consistency}},
\label{eq:joint_loss}
\end{equation}

where $\mathcal{L}_{\text{memory}}$ is token-level cross-entropy over
LLM-generated ground-truth summaries and
$\mathcal{L}_{\text{consistency}}$ penalizes divergence between the
action distribution conditioned on the predicted memory and that
conditioned on the ground-truth memory. We view training as the
next milestone once hot-memory integration is complete; the current
paper establishes the memory primitives required to make such
training meaningful.
